\clearpage
\begin{explanationbox}
  \textbf{\textcolor{CExplanation}{一、为什么会想到这种方法？}}

  \medskip
  先看一道不难的题：
  \[
    y=\frac{2x^2+3}{x^2+1}.
  \]
  它可以拆成 $y=2+\dfrac1{x^2+1}$，因此很容易看出 $2<y\le3$。
  可是换成
  \[
    y=\frac{x^2+1}{x^2+2x+2},
  \]
  分子、分母都带着 $x$，不容易像前一道题那样直接拆开。
  \textbf{困难不在于算出某个 $x$ 对应的 $y$，而在于找出所有可能的 $y$。}

  \medskip
  既然求的是 $y$ 的范围，不妨\textbf{暂时把 $y$ 当成一个已知数}，反过来解 $x$：
  \[
    \frac{x^2+1}{x^2+2x+2}=y
    \quad\Longleftrightarrow\quad
    (1-y)x^2-2yx+(1-2y)=0.
  \]
  如果这个方程找不到实数 $x$，说明原来的函数不可能等于这个 $y$。
  如果找得到，说明这个 $y$ 确实可以出现。
  这就把\textbf{求函数值域}变成了\textbf{判断方程有没有实数解}。

  \medskip
  \textbf{\textcolor{CExplanation}{二、为什么会出现判别式？}}

  \medskip
  初中学过，一元二次方程 $Ax^2+Bx+C=0$（$A\ne0$）的求根公式是
  \[
    x=\frac{-B\pm\sqrt{B^2-4AC}}{2A}.
  \]
  要算出实数 $x$，根号里的 $B^2-4AC$ 就不能小于 $0$。
  因此只需要计算一个关于 $y$ 的式子 $\Delta(y)$，再解 $\Delta(y)\ge0$。

  \medskip
  \textbf{\textcolor{CExplanation}{三、从图像上怎样理解？}}

  \medskip
  首页图中，蓝色边界水平线都能碰到曲线，橙色 $y=3$ 不能碰到。
水平线有交点才表示该输出值能取到；边界切触常对应 $\Delta=0$，但仍须确认二次项没有消失。

  \medskip
  \textbf{\textcolor{CExplanation}{四、两个值得记住的发现}}

  \medskip
  \textbf{发现 1：}若分母 $dx^2+ex+f$ 满足 $e^2-4df<0$，那么
  \[
    \Delta(y)=(e^2-4df)y^2+(4af+4cd-2be)y+(b^2-4ac)
  \]
  是一个关于 $y$ 的\textbf{开口向下的二次式}。所以解出的候选值总是一段连续的范围，而不是两头分开的范围。
  这段范围可能两端都能取到，也可能有一个端点取不到。

  \medskip
  \textbf{发现 2：}当 $x$ 的绝对值越来越大时，分子、分母主要由 $ax^2$ 与 $dx^2$ 决定，于是
  \[
    y\text{ 会越来越接近 }\frac ad.
  \]
  但是\textbf{越来越接近，不等于一定能够取到}。
  因此 $y=\dfrac ad$ 正好是我们必须另查的特殊值。

  \medskip
  \textbf{\textcolor{CExplanation}{五、什么时候优先用？}}
  分母始终不为 $0$、分子分母都是至多二次的式子，直接拆分又不明显时，可以优先把 $y$ 当成已知数，利用一元二次方程的实数解条件。
  如果能像第一道题那样很快拆成 $2+\dfrac1{x^2+1}$，就不必硬套判别式，选择更短的方法即可。
\end{explanationbox}
